\documentclass[11pt,oneside]{book}
\usepackage[margin=1.1in]{geometry}
\usepackage{xcolor}
\usepackage{booktabs}
\usepackage{tabularx}
\usepackage{titlesec}
\usepackage[hidelinks,bookmarksnumbered]{hyperref}

\definecolor{accent}{HTML}{5B2C6F}
\titleformat{\chapter}[display]
  {\normalfont\raggedright}
  {\color{accent}\Large\bfseries CHAPTER \thechapter}
  {0.5em}
  {\Huge\bfseries}
\titleformat{\section}{\Large\bfseries\color{accent}}{\thesection}{0.6em}{}

\newcommand{\publicationrecord}[6]{%
  \clearpage
  \thispagestyle{plain}
  {\color{accent}\rule{\textwidth}{1.2pt}}\\[1em]
  {\LARGE\bfseries Publication record\par}
  \vspace{1em}
  \begin{tabularx}{\textwidth}{@{}lX@{}}
  \toprule
  \textbf{Chapter} & #1 \\
  \textbf{Title} & #2 \\
  \textbf{Authors} & #3 \\
  \textbf{Status} & #4 \\
  \textbf{Citation or DOI} & #5 \\
  \bottomrule
  \end{tabularx}
  \vspace{1em}
  \section*{Candidate contribution}
  #6
  \section*{Rights and adaptation}
  Record the publishing agreement, reuse permission, version included, and all
  changes made for this thesis. Do not embed a publisher PDF unless the license
  or written permission explicitly allows it.
  \section*{Co-author confirmation}
  Candidate signature:\ \hrulefill\quad Date:\ \hrulefill\\[1.5em]
  Supervisor signature:\ \hrulefill\quad Date:\ \hrulefill
  \clearpage
}

\begin{document}

\begin{titlepage}
\centering
\vspace*{0.8in}
{\small THESIS INCLUDING PUBLISHED WORKS}\\[1em]
{\color{accent}\rule{0.8\textwidth}{1.5pt}}\\[1.2em]
{\Huge\bfseries Trustworthy Learning for\\
Longitudinal Health Data\par}
\vspace{1.2em}
{\Large Elena Marchetti}\\[0.6em]
{\large A thesis submitted for the degree of Doctor of Philosophy}

\vfill
{\large Aldermere University}\\
Department of Biomedical Informatics\\[0.5em]
July 2026
\vspace{0.7in}
\end{titlepage}

\frontmatter
\chapter*{Declaration of published and collaborative work}
This thesis contains three publication-format chapters. The table records each
work's status and the candidate's contribution. Full statements, permissions,
and co-author confirmations precede the corresponding chapters.

\begin{tabularx}{\textwidth}{@{}cXXl@{}}
\toprule
Chapter & Publication & Candidate role & Status \\
\midrule
2 & Missingness-aware clinical representations & Design, analysis, lead writing & Published \\
3 & Calibration across hospital systems & Experiments, software, lead writing & Accepted \\
4 & Prospective evaluation protocol & Design, protocol, joint writing & Submitted \\
\bottomrule
\end{tabularx}

\chapter*{Abstract}
Longitudinal health records are incomplete, irregular, and shaped by clinical
workflow. This thesis develops models that treat observation patterns as part
of the data-generating process. Three publication-format studies are connected
by a common methodological framework and an integrative discussion that
separates replicated findings from open questions.

\tableofcontents

\mainmatter
\chapter{Synthesis and research context}
A thesis by publication must make one coherent argument rather than merely bind
papers together. This chapter states the overarching question, relates the
methods across papers, resolves changes in terminology, and identifies work
performed specifically for the thesis.

\section{Integrated contribution}
The publications jointly show that missingness patterns improve prediction but
can also encode access disparities. The thesis contribution is a framework for
using those patterns while auditing transportability and subgroup calibration.

\publicationrecord
  {2}
  {Missingness-aware representations for clinical time series}
  {Elena Marchetti, Priya Nadeau, and Samuel Okoro}
  {Published}
  {Journal of Reliable Health AI 6 (2025), 101--124. DOI placeholder}
  {The candidate conceived the representation objective, implemented all
  experiments, analyzed the results, and drafted the manuscript. P.N. advised
  on study design. S.O. independently checked the statistical analysis.}

\chapter{Missingness-aware clinical representations}
Replace this chapter with the permitted manuscript version, reformatted to
match the thesis. Begin with a short bridge explaining how the paper advances
the thesis argument, and end with any thesis-specific limitations or updates.

\chapter{Integrative discussion}
Compare the evidence across papers, address contradictions, state the combined
limitations, and distinguish the thesis-level contribution from the
contribution of each publication.

\backmatter
\chapter*{Data, code, and publication rights}
List artifact identifiers and licenses, then retain permission records for each
included work. Confirm the exact institutional rules before submission.

\end{document}
